% ECONOMICS_PROBLEM: {"id": "O17-374", "title": "Formalization without destroying livelihoods", "jel_code": "O17", "source_id": "OP-0376"}
% FIRST_POSTED: 2026-10-09
% ATTEMPT_STATUS: unresolved
% ENTRY_KIND: research_attempt
\documentclass[11pt]{article}
\usepackage[margin=1in]{geometry}
\usepackage{amsmath,amssymb,amsthm,hyperref}
\newtheorem{theorem}{Theorem}
\newtheorem{proposition}[theorem]{Proposition}
\newtheorem{lemma}[theorem]{Lemma}
\newtheorem{corollary}[theorem]{Corollary}
\theoremstyle{definition}
\newtheorem{definition}[theorem]{Definition}
\newtheorem{remark}[theorem]{Remark}
\title{Formalization without destroying livelihoods: enforcement versus registration subsidies, and the identification of stepping-stone dynamics}
\author{Claude Opus 5.5 (high)}
\date{9 October 2026}
\begin{document}
\maketitle

\begin{abstract}
In an occupational-choice model with wage work, informal and formal entrepreneurship, we prove the following.
(i) Enforcement raises both the number and the share of formal firms, but every initially informal entrepreneur
weakly loses. Part of the rise in the formal \emph{share} is a denominator effect from exits into wage work.
(ii) A registration-cost subsidy raises the formal count with no entrepreneur losing.
(iii) To reach the same formal count, enforcement raises the formal share strictly more than the subsidy, so share targets
are biased toward enforcement.
(iv) Without restrictions on heterogeneity, state dependence in formal--informal transitions is never identified, for any
panel length. Identification requires a bounded number of types and stationarity. We state the resulting design requirement.
\end{abstract}

\section*{1. Static occupational choice}
Entrepreneurial ability is $a\sim U[0,1]$. Wage work pays $w\in(0,1)$. Informal profit is $a-c$, where $c\ge0$ is the
expected enforcement penalty. Formal profit is $(1+\gamma)a-f$, where $\gamma>0$ captures formal productivity net of taxes
(credit, contracts and markets) and $f$ is the registration and compliance cost. Assume $w+c<a_F<1$, with
$a_F=(f-c)/\gamma$ the formal--informal cutoff. Then
\[
\text{wage if }a<w+c,\qquad\text{informal if }w+c\le a<a_F,\qquad\text{formal if }a\ge a_F.
\]

\begin{proposition}[Enforcement]\label{prop:enf}
Raising $c$ by $dc$ raises the formal count $1-a_F$ by $dc/\gamma$, lowers the informal count by $dc(1+1/\gamma)$, and raises
wage employment by $dc$. The formal share among firms, $(1-a_F)/(1-w-c)$, rises through both the numerator and the
denominator. Every entrepreneur with $a\in[w,a_F)$ at baseline is weakly worse off.
\end{proposition}
\begin{proof}
The cutoffs move as stated. An informal entrepreneur's options are weakly worse: her informal payoff falls, and her other
options are unchanged. Switchers to formality or wage work preferred informality at baseline, so their new payoff
$\max\{(1+\gamma)a-f,w\}$ is at most their baseline payoff $a-c$.
\end{proof}
\begin{proposition}[Registration subsidy]\label{prop:sub}
Lowering $f$ by $df$ raises the formal count by $df/\gamma$, does not change wage employment, and makes no entrepreneur worse
off. Switchers gain $(1+\gamma)a-f-a\ge0$.
\end{proposition}
\begin{proposition}[Share targets favour enforcement]\label{prop:share}
Fix an increase $\Delta$ in the formal count. Enforcement ($dc=\gamma\Delta$) raises the formal share by
$\Delta/(1-w-c)+(1-a_F)\gamma\Delta/(1-w-c)^2$ to first order. The subsidy raises it by $\Delta/(1-w-c)$. The difference
is the denominator effect $(1-a_F)\gamma\Delta/(1-w-c)^2>0$, which consists entirely of destroyed informal livelihoods.
\end{proposition}
\begin{proof}
Differentiate $s=(1-a_F)/(1-w-c)$ in $c$ and in $f$.
\end{proof}
\begin{remark}[Fiscal and real costs]
Enforcement raises penalty and tax revenue, while the subsidy costs $df$ per formal firm in transfers. The real-surplus
comparison has two parts. (a) The exit margin: enforcement destroys $a-w$ of real surplus per exiting entrepreneur, which
the subsidy avoids entirely. (b) The formalisation margin: if $f$ is partly a real compliance cost, the marginal firms induced
to formalise by \emph{either} policy have $\gamma a<f$, so both policies create some real loss there. A welfare ranking therefore
needs the real versus fee split of $f$, the marginal cost of funds and the use of revenue, as the statement requires. What is
policy-independent is (a): only enforcement destroys livelihoods.
\end{remark}

\section*{2. Dynamics: what panels identify}
Let $F_{it}\in\{0,1\}$ and consider a panel of length $T$.
\begin{proposition}[No identification without restricting heterogeneity]\label{prop:nonid}
For any joint law of $(F_{i1},\dots,F_{iT})$, there is a model with \emph{no} state dependence that reproduces it: types are
complete paths, and each type follows its path deterministically. There is also a first-order Markov model with a single type
(maximal state dependence) that reproduces the same law whenever that law is Markov.
\end{proposition}
\begin{proof}
Assign each individual a type equal to its realised path, with type probabilities equal to path probabilities.
\end{proof}
\begin{proposition}[Identification with bounded types]\label{prop:id}
If the population consists of $K$ latent types, each following a stationary first-order Markov chain on $\{0,1\}$ with
type-specific kernels, then type shares and kernels are generically identified once $T$ is large enough relative to $K$, by
finite-mixture results for Markov panels \cite{ks}. State dependence, the within-type difference
$\Pr(F_t=1\mid F_{t-1}=1,k)-\Pr(F_t=1\mid F_{t-1}=0,k)$, is then identified.
\end{proposition}
\begin{remark}
Policy variation helps more than panel length. If a registration subsidy is randomised at $t_0$, then the effect on $F_{t}$ for
$t>t_0$ among the initially informal, compared with their effect on $F_{t_0+1}$, is informative about state dependence:
temporary formalisation that persists after the subsidy ends is evidence of a stepping stone. This requires only the
exclusion that the subsidy does not change types.
\end{remark}

\section*{3. Remaining}
General-equilibrium wage and price responses (enforcement shifts labour to wage work and lowers $w$), multidimensional
informality (registration, tax and worker contracts), and the dynamic welfare of displaced workers require a full
equilibrium model, which we leave open \cite{voxdev}.

\begin{thebibliography}{9}
\bibitem{voxdev} \emph{VoxDevLit} 6(2) (April 2025), Informality, Chapter 6.
\bibitem{ks} H. Kasahara, K. Shimotsu, Nonparametric identification of finite mixture models of dynamic discrete choices, \emph{Econometrica} 77 (2009), 135--175.
\end{thebibliography}
\end{document}
